% Programming with InitechBase, Chapter 14: Working with Two Databases
% (AN-16; annex Sec 8.2 P10).  Panels replayed (examples/p14-*.dot);
% sources cmd/flow.c SELECT, cmd/workarea.c (ten areas, each its own
% pointer), cmd/query.c SEEK; a field of another area is reached through a
% memory variable (the shipped evaluator resolves names in the selected
% area).
\chapter{Working with Two Databases}
\Gates{test-pubs-examples, test-interp-use, test-samir-query}
The departments of Initech are kept in their own database, \file{DEPTS.DBF},
with the name and the floor of each; the staff roster gives only the code.  A
program that prints an employee with the name of the department looks the code
up in the second database.  This chapter shows the way InitechBase programs do
it: two work areas, a memory variable, and \cmd{SEEK}.

\begin{ChapterCovers}
\item Two databases in two work areas
\item Carrying a value from one to the other
\item Looking up a code with \cmd{SEEK}
\item Counting the records of another database
\end{ChapterCovers}

\ATEHead{Preparing for This Chapter}
Start InitechBase with a fresh copy of the sample diskette.

\ATEHead{Two Work Areas}
Each database is opened in its own work area, at the dot prompt, before the
program runs.  The program selects the work area of the database it wants to
read; each work area keeps its own record pointer, so moving in one does not
move the other.

A field name always refers to the database of the selected work area.  To use
a value of the staff roster while the departments are selected, the program
stores it in a memory variable first.

\ATEHead{Looking Up the Department}
\file{STAFFDEP.PRG} lists every employee with the name of the department:
\ProgramListing{STAFFDEP.PRG}
\begin{Steps}
\item Type \type{SELECT 2}, then \type{USE DEPTS INDEX DEPTS}.
\item Type \type{SELECT 1}, then \type{USE STAFF}.
\item Type \type{DO STAFFDEP}.
\end{Steps}
\DotExample{examples/p14-staffdep}
For each employee the program keeps the code in \texttt{MDEPT}, selects the
departments, seeks the code in their index and takes the name; then it selects
the staff again before \cmd{SKIP}.  A code with no department would print the
words in parentheses instead of a name.

\ATEHead{Counting Through an Index}
\file{HEADCNT.PRG} works the other way round: for each department it counts the
employees, through the index of the staff on the department:
\ProgramListing{HEADCNT.PRG}
\DotExample{examples/p14-headcnt}
The departments are read in their order on the disk; for each, the inner loop
seeks the first employee of the department and counts while the department
stays the same.

\ATEHead{Where Do You Go From Here?}
Chapter~15 prints a long listing in pages.

\ATEHead{Summary}
\begin{enumerate}
\item Open each database in its own work area before the program runs.
\item \cmd{SELECT} chooses the database whose fields a name refers to.
\item Carry a value from one database to another in a memory variable.
\item \cmd{SEEK} looks the value up in the index of the other database.
\end{enumerate}
